\section{Logging and Debugging Tools}
\label{sec:implementation:logging-debugging}

Comprehensive logging is essential for safety validation, operational transparency, and post-event analysis in an interlocking system. The proposed framework adopts a multi-tiered approach that captures events with varying levels of detail, ensuring that both real-time monitoring and long-term auditing needs are met. At the backend level, every interlocking decision—such as a signal clearance or a rejected route request—is timestamped with its input parameters and validation outcomes, guaranteeing complete traceability of operational behavior. Structured logs are stored in dedicated database tables, where fields capture component identity, action taken, status, and any associated failure reason. This structured format allows efficient querying for audits, diagnostics, and compliance checks.

During development and simulation, the system leverages Qt’s console utilities to provide categorized diagnostic messages. Routine updates are emitted as standard debug logs, warnings highlight unexpected but non-critical anomalies, and critical errors are flagged immediately to indicate failed safety validations or major exceptions. This real-time feedback accelerates troubleshooting by giving engineers immediate insight into system behavior under different conditions.

Beyond live debugging, the logging infrastructure is equally important for post-mortem analysis. By reconstructing event sequences from historical logs, engineers can validate whether safety logic responded correctly under stress or identify where faults originated. This dual role—supporting both immediate debugging and long-term accountability—ensures that the interlocking system remains transparent, verifiable, and aligned with the rigorous demands of railway safety.